% Part: normal-modal-logic
% Chapter: frame-definability
% Section: definability

\documentclass[../../../include/open-logic-section]{subfiles}

\begin{document}

\olfileid{nml}{frd}{def}

\olsection{কাঠামো-সংজ্ঞায়নযোগ্যতা}

\olref[acc]{thm:soundschemas}-এর বিপরীত নিহিতার্থগুলি মডেলের জন্য ব্যর্থ হলেও, ``মডেল''-এর জায়গায় ``কাঠামো'' বসালে সেগুলি সত্য হয়: \olref[acc]{thm:soundschemas}-এ আলোচিত ধর্মগুলির ক্ষেত্রে, কোনো \emph{কাঠামোতে} !!a{formula} বৈধ হলে সেই কাঠামোর অভিগম্যতা সম্পর্কের সংশ্লিষ্ট ধর্মটি \emph{থাকবেই}। তাই আলোচিত !!{formula}s সংশ্লিষ্ট ধর্মবিশিষ্ট কাঠামোর শ্রেণিকে \emph{সংজ্ঞায়িত করে}।

\begin{defn}
  $\mClass{F}$ কাঠামোর একটি শ্রেণি হলে, $\mModel{F} \Entails !A$ ঠিক সেই সব কাঠামো~$\mModel{F} \in \mClass{F}$-এর ক্ষেত্রেই সত্য হলে আমরা বলি, $!A$ শ্রেণি~$\mClass{F}$-কে \emph{সংজ্ঞায়িত করে}।
\end{defn}

এবার কাঠামোর জন্য সংজ্ঞায়নযোগ্যতার পূর্ণ ফলগুলি প্রতিষ্ঠা করা যাক।

\begin{thm}\ollabel{thm:fullCorrespondence}
\olref[acc]{tab:five}-এর ডান দিকের !!{formula} কোনো কাঠামো~$\mModel{F}$-এ বৈধ হলে, $\mModel{F}$-এর বাম দিকের ধর্মটি থাকে।
\end{thm}

\begin{proof}
  \begin{enumerate}
  \item ধরা যাক, $\mModel{F} = \tuple{W, R}$-এ \Ax{D} বৈধ, অর্থাৎ $\mModel{F} \Entails \Box p \lif \Diamond p$। $\mModel{F}$-এর উপর প্রতিষ্ঠিত একটি মডেল $\mModel{M} =
    \tuple{W, R, V}$ এবং $w \in
    W$ নিই। দেখাতে হবে, কোনো $v$ আছে যার জন্য~$Rwv$। বিপরীতটি ধরা যাক: তাহলে~$!A$ যাই হোক না কেন, $\mSat{M}{\Box !A}[w]$ এবং $\mSat/{M}{\Diamond !A}[w]$ উভয়ই সত্য; বিশেষত~$p$-এর ক্ষেত্রেও। কিন্তু তখন $\mSat/{M}{\Box p \lif
      \Diamond p}[w]$, যা $\mModel{F}
    \Entails \Box p \lif \Diamond p$ অনুমানের বিরোধী।
    % BN-SRC-343: উৎসে প্রথম সত্যতা-দাবির জগৎ-সূচক বাদ পড়েছে।
  \item ধরা যাক, $\mModel{F}$-এ \Ax{T} বৈধ, অর্থাৎ $\mModel{F} \Entails \Box
    p \lif p$। যেকোনো জগৎ $w \in W$ নিই; দেখাতে হবে $Rww$। $u \in V(p)$ তখনই এবং কেবল তখনই ধরি, যখন $Rwu$ ($p$ ছাড়া অন্য $q$-এর জন্য $V(q)$ নির্বিচার; যেমন $V(q) = \emptyset$)। ধরা যাক,
    $\mModel{M} = \tuple{W, R, V}$। নির্মাণ অনুযায়ী, $Rwu$ হয় এমন সব $u$-এর জন্য $\mSat{M}{p}[u]$; তাই $\mSat{M}{\Box p}[w]$। কিন্তু অনুমান অনুযায়ী $w$-তে $\Box p \lif p$ সত্য। সুতরাং
    $\mSat{M}{p}[w]$। অথচ $V$-এর সংজ্ঞা অনুযায়ী এটি কেবল~$Rww$ হলেই সম্ভব।
    % BN-SRC-344: মানায়নের পরের বন্ধনীটি গাণিতিক সীমার বাইরে রাখা হয়েছে।
  \item বিপরীত-অনুবর্তন প্রমাণ করি: ধরা যাক, $\mModel{F}$ প্রতিসম নয়। দেখাব, \Ax{B}, অর্থাৎ $p \lif \Box\Diamond p$, $\mModel{F}= \tuple{W, R}$-এ বৈধ নয়। $\mModel{F}$ প্রতিসম না হলে এমন $u$, $v \in W$ আছে যে $Ruv$ কিন্তু $Rvu$ নয়। $V$ এমনভাবে সংজ্ঞায়িত করি, যেন $w \in V(p)$ তখনই এবং কেবল তখনই, যখন $Rvw$ নয় (অন্যান্য ক্ষেত্রে $V$ নির্বিচার)। ধরা যাক, $\mModel{M} =\tuple{W, R,
      V}$। এখন $V$-এর সংজ্ঞা অনুযায়ী, $Rvw$ নয় এমন সব $w$-এর জন্য $\mSat{M}{p}[w]$; বিশেষত, $Rvu$ নয় বলে $\mSat{M}{p}[u]$। আবার $Rvw$ তখনই এবং কেবল তখনই, যখন $w \notin V(p)$। তাই এমন কোনো~$w$ নেই যার জন্য $Rvw$ এবং $\mSat{M}{p}[w]$; সুতরাং $\mSat/{M}{\Diamond p}[v]$। যেহেতু~$Ruv$, তাই $\mSat/{M}{\Box\Diamond p}[u]$-ও। ফলে $\mSat/{M}{p \lif
      \Box\Diamond p}[u]$; অতএব $\Ax{B}$~$\mModel{F}$-এ বৈধ নয়।
  \item ধরা যাক, $\mModel{F} = \tuple{W,R}$-এ \Ax{4} বৈধ, অর্থাৎ $\mModel{F} \Entails \Box p \lif \Box\Box p$। আবার $u$, $v$, $w \in W$ এমন নির্বিচার জগৎ যে $Ruv$ এবং $Rvw$; দেখাতে হবে $Ruw$। $V$ এমনভাবে সংজ্ঞায়িত করি, যেন $z \in V(p)$ তখনই এবং কেবল তখনই, যখন $Ruz$ (অন্যত্র $V$ নির্বিচার)। ধরা যাক, $\mModel{M} =\tuple{W, R, V}$। $V$-এর সংজ্ঞা অনুযায়ী $Ruz$ হয় এমন সব $z$-এর জন্য $\mSat{M}{p}[z]$; তাই $\mSat{M}{\Box p}[u]$। কিন্তু অনুমান অনুযায়ী \Ax{4}, অর্থাৎ $\Box p
    \lif \Box \Box p$, $u$-তে সত্য। ফলে $\mSat{M}{\Box \Box
      p}[u]$। $Ruv$ ও $Rvw$ বলে $\mSat{M}{p}[w]$; আর~$V$-এর সংজ্ঞা অনুযায়ী এটি কেবল $Ruw$ হলেই সম্ভব, যা দেখানোর ছিল।
  \item বিপরীত-অনুবর্তনে এগোই: কাঠামো $\mModel{F} = \tuple{W, R}$ ইউক্লিডীয় নয় ধরে দেখাব, এতে~\Ax{5} ব্যর্থ হয়; অর্থাৎ $\mModel{F} \Entails/ \Diamond p \lif
      \Box\Diamond p$। ধরা যাক, $u$, $v$, $w \in W$ এমন জগৎ যে $Rwu$ ও $Rwv$, কিন্তু~$Ruv$ নয়। $V$ এমনভাবে সংজ্ঞায়িত করি, যেন সব জগৎ $z$-এর জন্য $z \in V(p)$ তখনই এবং কেবল তখনই, যখন~$Ruz$ \emph{নয়}। ধরা যাক, $\mModel{M} =\tuple{W, R, V}$। তাহলে অনুমান থেকে $\mSat{M}{p}[v]$ এবং $Rwv$ বলে $\mSat{M}{\Diamond p}[w]$। কিন্তু $Ruy$ ও $\mSat{M}{p}[y]$ উভয়ই সত্য এমন কোনো জগৎ $y$ নেই। তাই $\mSat/{M}{\Diamond
      p}[u]$। $Rwu$ থেকে পাই $\mSat/{M}{\Box\Diamond
      p}[w]$; অতএব \Ax{5}, অর্থাৎ $\Diamond p \lif \Box\Diamond p$, জগৎ~$w$-তে ব্যর্থ।
  \end{enumerate}
\end{proof}

\Ax{D}-এর প্রমাণ ও অন্য ক্ষেত্রগুলির প্রমাণের মধ্যে একটি পার্থক্য লক্ষ করুন: মানায়ন~$V$-এর উল্লেখ করতে হয়নি। প্রকৃতপক্ষে আমরা প্রমাণ করেছি, $\mSat{M}{\Ax{D}}$ হলে $\mModel{M}$ সিরিয়াল। তাই \Ax{D} শুধু কাঠামোর নয়, সিরিয়াল \emph{মডেলের} শ্রেণিও সংজ্ঞায়িত করে।

\begin{cor}\ollabel{prop:D-serial}
  যে মডেলে \Ax{D} সত্য, সেটি সিরিয়াল।
\end{cor}

\begin{cor}
\olref[acc]{tab:five}-এর ডান দিকের প্রত্যেক !!{formula} বাম দিকের সংশ্লিষ্ট ধর্মবিশিষ্ট কাঠামোর শ্রেণিকে সংজ্ঞায়িত করে।
\end{cor}

\begin{proof}
  \olref[acc]{thm:soundschemas}-এ প্রমাণ করেছি, কোনো মডেলের বাম দিকের ধর্মটি থাকলে ডান দিকের !!{formula} তাতে সত্য। ফলে কাঠামো~$\mModel{F}$-এর বাম দিকের ধর্মটি থাকলে ডান দিকের !!{formula}~$\mModel{F}$-এ বৈধ। \olref{thm:fullCorrespondence}-এ বিপরীত নিহিতার্থগুলি প্রমাণ করেছি: ডান দিকের !!a{formula}~$\mModel{F}$-এ বৈধ হলে, $\mModel{F}$-এর বাম দিকের ধর্মটি থাকে।
\end{proof}

\begin{prob}
দেখান, \olref[nml][frd][acc]{tab:anotherfive}-এর ডান দিকের !!{formula} কোনো কাঠামো~$\mModel{F}$-এ বৈধ হলে, $\mModel{F}$-এর বাম দিকের ধর্মটি থাকে। এর জন্য বাম দিকের ধর্মটি \emph{পূরণ করে না} এমন কাঠামো নিয়ে উপযুক্ত~$V$ সংজ্ঞায়িত করুন, যাতে ডান দিকের !!{formula} কোনো জগতে মিথ্যা হয়।
\end{prob}

\olref{thm:fullCorrespondence} আরও দেখায় যে ধর্মগুলি একত্র করা যায়: যেমন,~$\mModel{F}$-এ \Ax{B} ও \Ax{4} উভয়ই বৈধ হলে কাঠামোটি প্রতিসম ও পরিযায়ী—উভয়ই; ইত্যাদি। বহু গুরুত্বপূর্ণ মোডাল যুক্তিকে কয়েকটি কাঠামোগত ধর্ম সমন্বিত সব কাঠামোতে বৈধ !!{formula}s-এর সমষ্টি হিসেবে চিহ্নিত করা হয়। তাই সংশ্লিষ্ট সংজ্ঞায়ক !!{formula}s বৈধ এমন সব কাঠামোতে বৈধ !!{formula}s-এর সমষ্টি হিসেবেও সেগুলিকে চিহ্নিত করা যায়। যেমন, প্রচলিত তন্ত্র \Log{S4} হলো সব প্রতিবিম্বী ও পরিযায়ী কাঠামোতে বৈধ !!{formula}s-এর সমষ্টি; অর্থাৎ, যে কাঠামোগুলিতে \Ax{T} ও~\Ax{4} উভয়ই বৈধ। \Log{S5} হলো সব প্রতিবিম্বী, প্রতিসম ও ইউক্লিডীয় কাঠামোতে বৈধ সূত্রের সমষ্টি; অর্থাৎ, যে কাঠামোগুলিতে \Ax{T}, \Ax{B} এবং \Ax{5} সবই বৈধ।

$R$-এর ধর্মগুলির মধ্যকার যৌক্তিক সম্পর্ক সাধারণত তাদের সংজ্ঞায়ক !!{formula}s-এর সম্পর্কের সঙ্গে মেলে। যেমন, প্রতিটি প্রতিবিম্বী সম্পর্ক সিরিয়াল; তাই কোনো কাঠামোতে \Ax{T} বৈধ হলে~\Ax{D}-ও বৈধ। (লক্ষ করুন, এটি \emph{অনুবর্তনের} সম্পর্ক নয়। $\mSat{M}{\Ax{T}}[w]$ হলেই $\mSat{M}{\Ax{D}}[w]$ হবে, তা নয়।) এ রকম কয়েকটি সম্পর্ক নথিবদ্ধ করি।

\begin{prop}\ollabel{prop:relation-facts}
  ধরা যাক, $R$ সমষ্টি $W$-এর উপর একটি দ্বিপদী সম্পর্ক। তাহলে:
  \begin{enumerate}
  \item $R$ প্রতিবিম্বী হলে সিরিয়াল।
  \item $R$ প্রতিসম হলে, সেটি পরিযায়ী তখনই এবং কেবল তখনই, যখন সেটি ইউক্লিডীয়।
  \item $R$ প্রতিসম বা ইউক্লিডীয় হলে সেটি দুর্বলভাবে অভিমুখী (এর ``হীরক ধর্ম'' আছে)।
  \item $R$ ইউক্লিডীয় হলে সেটি দুর্বলভাবে সংযুক্ত।
  \item $R$ অপেক্ষকধর্মী হলে সেটি সিরিয়াল।
  \end{enumerate}
\end{prop}

\begin{prob}
  \olref[nml][frd][def]{prop:relation-facts} প্রমাণ করুন।
\end{prob}

\end{document}
